\documentclass[a4paper,11pt]{amsart}

% --- IDIOMA Y CODIFICACIÓN ------------------------------------------------
\usepackage[utf8]{inputenc}
\usepackage[T1]{fontenc}
\usepackage[spanish,es-tabla,es-nodecimaldot]{babel}

% --- TIPOGRAFÍA (STIX Two: texto y matemáticas, la de muchas revistas) -------
\usepackage{stix2}
\linespread{1.07}
\usepackage[final]{microtype}

% --- MATEMÁTICAS Y UNIDADES ------------------------------------------------
\usepackage{amsmath, mathtools}
\usepackage{bm}
\usepackage{siunitx}
\sisetup{
  output-decimal-marker = {.},
  group-minimum-digits  = 4,
  group-digits          = integer,
  group-separator       = {\,},
  per-mode              = symbol,
  separate-uncertainty  = true,
  range-phrase          = {\text{ a }},
  list-final-separator  = {\text{ y }},
  list-pair-separator   = {\text{ y }},
  detect-all            = false
}

% --- PÁGINA, COLORES Y GRÁFICOS --------------------------------------------
\usepackage[a4paper,margin=2.6cm,footskip=1.45cm]{geometry}
\usepackage{graphicx}
\usepackage{xcolor}
% Gama de azules: del marino al aguamarina, con el turquesa como acento
\definecolor{azulprof}{HTML}{0B1F3A}     % azul marino profundo
\definecolor{azulacero}{HTML}{1F5A96}    % azul cobalto
\definecolor{azulpetroleo}{HTML}{12667A} % azul petróleo
\definecolor{turquesa}{HTML}{17A2A6}     % turquesa
\definecolor{aguamarina}{HTML}{86D3CF}   % aguamarina
\definecolor{azulclaro}{HTML}{B8DCEF}    % azul cielo pálido
\definecolor{grispizarra}{HTML}{44556B}  % gris pizarra azulado
\definecolor{grisclaro}{HTML}{E2EBF3}    % gris perla azulado
\definecolor{azulfondo}{HTML}{EFF6FB}    % fondo tenue de referencias y resultados
\colorlet{acento}{turquesa!85!azulprof}  % curvas teóricas y comparaciones
\usepackage{tikz}
\usetikzlibrary{babel, arrows.meta, calc, positioning, decorations.pathmorphing,
                decorations.pathreplacing, decorations.markings, patterns, shadings}
\usepackage{pgfplots}
\usepgfplotslibrary{groupplots, fillbetween}
\pgfplotsset{compat=1.17}
\usepackage{booktabs}
\usepackage[section]{placeins}
\usepackage{multirow}
\usepackage[most]{tcolorbox}
\usepackage{array}
\usepackage{enumitem}
\setlist[itemize]{leftmargin=1.6em, itemsep=2pt, topsep=4pt}
\setlist[enumerate]{leftmargin=1.8em, itemsep=2pt, topsep=4pt}

% --- ENLACES Y REFERENCIAS CRUZADAS ----------------------------------------
\usepackage[colorlinks=true, linkcolor=azulacero, citecolor=azulacero,
            urlcolor=azulacero, plainpages=false, pdfusetitle=false,
            pdftitle={Práctica 4: Péndulo de Pohl},
            pdfauthor={David Gómez García; Daniel Martín Guerrero}]{hyperref}
\usepackage[spanish,nameinlink]{cleveref}

% --- BIBLIOGRAFÍA -----------------------------------------------------------
\usepackage{csquotes}
\usepackage[backend=biber, style=numeric, sorting=none, giveninits=true,
            maxbibnames=6, minbibnames=3, doi=true, url=true,
            isbn=false]{biblatex}
\addbibresource{referencias.bib}
\DeclareFieldFormat[article,techreport,unpublished,misc]{title}{#1\isdot}
\DeclareFieldFormat{url}{\url{#1}}
\renewcommand*{\bibfont}{\small}
\setlength{\bibitemsep}{3pt}

% --- PIEZAS DE DIBUJO DEL PÉNDULO DE POHL ------------------------------------
%%%%%%%%%%%%%%%%%%%%%%%%%%%%%%%%%%%%%%%%%%%%%%%%%%%%%%%%%%%%%%%%%%%%%%%%
% Piezas de dibujo del péndulo de Pohl (TikZ)
%
% Vista frontal y minimalista del aparato del laboratorio: la escala
% graduada fija, el anillo de cobre con sus radios, el muelle en espiral
% del eje, el fiel, el electroimán que abraza la parte inferior del anillo
% y el excitador (palanca, biela y disco excéntrico del motor).
% Se usan en la portada y en el esquema del montaje experimental.
% Requiere los colores azulprof, azulacero, azulpetroleo, turquesa,
% aguamarina, azulclaro, grispizarra y grisclaro (definidos en main.tex).
%%%%%%%%%%%%%%%%%%%%%%%%%%%%%%%%%%%%%%%%%%%%%%%%%%%%%%%%%%%%%%%%%%%%%%%%

% Escala graduada fija: corona con ranuras cada 6 grados (más largas cada 30)
% \escalapohl{radio interior}{radio exterior}
\newcommand{\escalapohl}[2]{%
  \fill[azulprof, even odd rule] (0,0) circle (#2) (0,0) circle (#1);
  \foreach \a in {0,6,...,354}{%
    \pgfmathtruncatemacro{\largo}{mod(\a,30)==0 ? 1 : 0}%
    \ifnum\largo=1
      \draw[white, line width=0.9pt, line cap=round]
        (\a:{#1+0.34*(#2-#1)}) -- (\a:{#2-0.16*(#2-#1)});
    \else
      \draw[white!80!azulclaro, line width=0.45pt, line cap=round]
        (\a:{#1+0.52*(#2-#1)}) -- (\a:{#2-0.18*(#2-#1)});
    \fi}}

% Anillo de cobre (gira con el ángulo #1): corona, cuatro radios y cubo
% \anillopohl{ángulo}{radio interior}{radio exterior}
\newcommand{\anillopohl}[3]{%
  \begin{scope}[rotate=#1]
    \foreach \a in {45,135,225,315}{%
      \fill[azulpetroleo!85!azulprof, rounded corners=0.4pt]
        (\a-4:0.32) -- (\a-1.6:#2+0.05) -- (\a+1.6:#2+0.05) -- (\a+4:0.32) -- cycle;}
    \fill[turquesa!78!azulpetroleo, even odd rule] (0,0) circle (#3) (0,0) circle (#2);
    \draw[azulpetroleo!80!azulprof, line width=0.5pt] (0,0) circle (#3);
    \draw[azulpetroleo!80!azulprof, line width=0.5pt] (0,0) circle (#2);
    \draw[white, opacity=0.35, line width=0.6pt] (100:{#3-0.07}) arc[start angle=100, end angle=165, radius={#3-0.07}];
    % fiel: aguja solidaria con el anillo que apunta a la escala
    \fill[white, draw=azulprof, line width=0.45pt]
      (90:{#2-0.1}) -- ($(90:{#3+0.12})+(-0.075,0)$) -- (90:{#3+0.52})
      -- ($(90:{#3+0.12})+(0.075,0)$) -- cycle;
  \end{scope}}

% Muelle en espiral del eje (espiral de Arquímedes); su extremo exterior
% gira con el excitador (ángulo #1)
% \muelleespiral{ángulo del extremo exterior}{radio interior}{radio exterior}{vueltas}
\newcommand{\muelleespiral}[4]{%
  \draw[azulprof, line width=0.85pt, line cap=round, domain=0:#4*360, samples=220, smooth, variable=\t]
    plot ({(#2+(#3-#2)*\t/(#4*360))*cos(\t+#1)}, {(#2+(#3-#2)*\t/(#4*360))*sin(\t+#1)});
  \fill[grisclaro, draw=azulprof, line width=0.5pt] (0,0) circle (#2);
  \fill[azulprof] (0,0) circle (0.045);}

% Electroimán: dos bobinas que abrazan la parte inferior del anillo
% \electroiman{radio medio del anillo}{activo: 1 dibuja el campo}
\newcommand{\electroiman}[2]{%
  \begin{scope}[shift={(0,-#1)}]
    \fill[grisclaro, draw=grispizarra, line width=0.45pt, rounded corners=1pt]
      (-0.62,-0.46) rectangle (0.62,0.30);
    \foreach \x in {-0.48,-0.36,...,0.5}{%
      \draw[azulacero!55, line width=0.35pt] (\x,-0.40) -- (\x,0.24);}
    \fill[azulprof, rounded corners=0.8pt] (-0.72,-0.62) rectangle (0.72,-0.46);
    \ifnum#2=1
      \foreach \x in {-0.3,0.3}{%
        \draw[azulprof, line width=0.45pt, fill=white] (\x,-0.08) circle (0.11);
        \draw[azulprof, line width=0.45pt] (\x,-0.08) ++(-0.075,-0.075) -- ++(0.15,0.15)
          (\x,-0.08) ++(-0.075,0.075) -- ++(0.15,-0.15);}
    \fi
  \end{scope}}

% Excitador: palanca solidaria con el extremo exterior del muelle, biela
% y disco excéntrico del motor (fase de giro #1, en grados)
% \excitador{fase}{punto de la palanca}{centro del disco}{radio de la excéntrica}
\newcommand{\excitador}[4]{%
  \coordinate (PinE) at ($(#3)+(#1:#4)$);
  % motor: carcasa discreta y disco
  \fill[grisclaro, draw=grispizarra!80, line width=0.4pt, rounded corners=2pt]
    ($(#3)+(0.25,-0.62)$) rectangle ($(#3)+(2.05,0.55)$);
  \foreach \y in {0.30,0.18,0.06}{\draw[grispizarra!45, line width=0.4pt]
    ($(#3)+(0.55,\y)$) -- ($(#3)+(1.0,\y)$);}
  \fill[grispizarra!80] ($(#3)+(1.45,-0.08)$) circle (0.09);
  \fill[grispizarra!80] ($(#3)+(1.72,-0.32)$) circle (0.08);
  \fill[azulprof] (#3) circle (0.62);
  \draw[white, line width=0.5pt] (#3) circle (0.48);
  \fill[white] (#3) circle (0.05);
  % biela
  \draw[azulacero, line width=1.3pt, line cap=round] (#2) -- (PinE);
  \fill[white, draw=azulprof, line width=0.5pt] (#2) circle (0.075);
  \fill[white, draw=azulprof, line width=0.5pt] (PinE) circle (0.075);}

% Péndulo de Pohl completo
% \pohl{ángulo del anillo}{fase del motor}{dibuja el campo (0/1)}
\newcommand{\pohl}[3]{%
  \coordinate (Palanca) at (-1.15,-3.95);
  \coordinate (Motor) at (4.05,-4.25);
  \escalapohl{2.85}{3.45}
  \anillopohl{#1}{2.02}{2.42}
  \electroiman{2.22}{#3}
  % palanca del excitador: arranca del extremo exterior del muelle
  \draw[azulprof, line width=2.1pt, line cap=round] (-106:1.1) -- (Palanca);
  \muelleespiral{110}{0.17}{1.1}{3.4}
  \excitador{#2}{Palanca}{Motor}{0.36}}

% Ornamento: una oscilación amortiguada con su envolvente
% \ondita[escala]
\newcommand{\ondita}[1][1]{%
  \tikz[baseline=-0.55ex, x=#1*1pt, y=#1*1pt]{%
    \draw[azulprof!55, line width=0.3pt, densely dotted]
      plot[domain=0:18, samples=30, smooth, variable=\t] ({\t},{2.8*exp(-\t/9)});
    \draw[turquesa, line width=0.45pt, line join=round]
      plot[domain=0:18, samples=90, smooth, variable=\t]
      ({\t},{2.8*exp(-\t/9)*cos(1.9*\t r)});}}


% =========================================================================
%        CABECERAS, PIE DE PÁGINA, SECCIONES E ÍNDICE
% =========================================================================
% Ornamento entre secciones: una oscilación amortiguada entre dos filetes
\newcommand{\ornamento}{%
  \par\vspace{10pt}{\centering
  \tikz{\draw[azulclaro, line width=0.4pt] (-1.6,0) -- (-0.4,0);
        \draw[azulclaro, line width=0.4pt] (0.4,0) -- (1.6,0);
        \node at (0,0) {\ondita[1.25]};}\par}\vspace{6pt}}

\makeatletter
\def\ps@headings{%
  \let\@mkboth\markboth
  \def\@evenhead{\normalfont\footnotesize\color{grispizarra}\scshape
    \textls[60]{Laboratorio de Mecánica y Ondas}\hfil
    \textls[60]{Péndulo de Pohl\,\textperiodcentered\,Oscilaciones forzadas}}%
  \let\@oddhead\@evenhead
  \def\@evenfoot{\hfil\normalfont\small\color{azulacero}%
    \ondita\enspace\thepage\enspace\ondita\hfil}%
  \let\@oddfoot\@evenfoot}
\def\ps@firstpage{\ps@plain\def\@oddhead{}\let\@evenhead\@oddhead
  \def\@oddfoot{\hfil\normalfont\small\color{azulacero}%
    \ondita\enspace\thepage\enspace\ondita\hfil}%
  \let\@evenfoot\@oddfoot}
\pagestyle{headings}

% Secciones: número en azul cobalto, título en versalitas espaciadas; debajo,
% un subtítulo en cursiva y un filete rematado en un pequeño círculo turquesa.
\def\section{\@startsection{section}{1}%
  \z@{1.3\linespacing\@plus.8\linespacing}{.25\linespacing}%
  {\normalfont\Large\color{azulprof}\scshape\lsstyle\raggedright}}
\let\@seccntformatorig\@seccntformat
\def\@seccntformat#1{%
  \ifnum\pdfstrcmp{#1}{section}=\z@
    {\color{azulacero}\upshape\csname the#1\endcsname}\hspace{0.55em}%
  \else\@seccntformatorig{#1}\fi}
% Subsecciones: en negrita azul marino, con el número en azul cobalto
\def\subsection{\@startsection{subsection}{2}%
  \z@{.9\linespacing\@plus.6\linespacing}{.3\linespacing}%
  {\normalfont\bfseries\color{azulprof}}}
\def\subsubsection{\@startsection{subsubsection}{3}%
  \z@{.7\linespacing\@plus.5\linespacing}{.2\linespacing}%
  {\normalfont\itshape\color{azulacero}}}
\newcommand{\filete}{%
  \tikz[baseline]{\draw[azulclaro, line width=0.45pt] (0,0) -- (\linewidth-5pt,0);
    \draw[turquesa, line width=0.55pt, fill=white] (\linewidth-2.6pt,0) circle (2.3pt);}}
\DeclareRobustCommand{\tocsub}[1]{%
  \ \textperiodcentered\ {\itshape\color{grispizarra}#1}}
\newcommand{\subtitulo}[1]{%
  \par\nopagebreak\vspace{-2pt}%
  \noindent{\itshape\color{grispizarra}#1}\par\nopagebreak\vspace{1pt}%
  \noindent\filete\par\nopagebreak\vspace{0.55\baselineskip}}
% \seccion{Título}{Subtítulo}: sección numerada con subtítulo (índice incluido).
\let\@tocextra\@empty
\def\@tocwriteb#1#2#3{%
  \begingroup
    \def\@tocline##1##2##3##4##5##6{%
      \ifnum##1>\c@tocdepth
      \else \sbox\z@{##5\let\indentlabel\@tochangmeasure##6}\fi}%
    \csname l@#2\endcsname{#1{\csname#2name\endcsname}{\@secnumber}{}}%
  \endgroup
  \addcontentsline{toc}{#2}%
    {\protect#1{\csname#2name\endcsname}{\@secnumber}{#3\@tocextra}}%
  \global\let\@tocextra\@empty}
\newcommand{\seccion}[2]{%
  \gdef\@tocextra{\tocsub{#2}}%
  \section{#1}\subtitulo{#2}}

% Índice: número en azul cobalto, título en azul marino y subtítulo en cursiva
\def\l@section{\@tocline{1}{7pt}{0pt}{}{}}
\renewcommand{\tocsection}[3]{%
  \indentlabel{\@ifnotempty{#2}{{\color{azulacero}#2}\quad}}{\color{azulprof}#3}}
\def\l@subsection{\@tocline{2}{1pt}{1.55pc}{2.4pc}{\small}}
\renewcommand{\contentsnamefont}{\scshape\color{azulprof}\lsstyle}
\makeatother
\numberwithin{equation}{section}
% Párrafos: encabezado en cursiva, en línea con el texto
\makeatletter
\def\paragraph{\@startsection{paragraph}{4}{\z@}{6pt plus 2pt}{-0.6em}%
  {\normalfont\itshape\color{azulprof}}}
\makeatother
\setcounter{tocdepth}{2}
\makeatletter
\def\@@and{y}
\renewcommand{\abstractname}{\color{azulacero}\lsstyle Resumen}
\makeatother

% =========================================================================
%               NOTACIÓN Y MACROS DE USO FRECUENTE
% =========================================================================
\newcommand{\dif}{\mathrm{d}}
\newcommand{\ii}{\mathrm{i}}                  % unidad imaginaria
\newcommand{\wz}{\omega_0}                    % frecuencia natural (sin amortiguamiento)
\newcommand{\wa}{\omega}                      % frecuencia de la oscilación amortiguada
\newcommand{\WR}{\Omega_{\mathrm{R}}}         % frecuencia de resonancia
\newcommand{\bR}{\beta_{\mathrm{R}}}          % amplitud de resonancia
\newcommand{\Tcoef}{t_{95,\nu}}               % factor de cobertura al 95 %
\newcommand{\ext}[1]{U(#1)}                   % incertidumbre ampliada
\newcommand{\mean}[1]{\overline{#1}}
% Dato pendiente de incorporar (se resalta con discreción)
\newcommand{\pend}[1]{\textcolor{grispizarra}{\colorbox{grisclaro}{\strut\,#1\,}}}
\newcommand{\pendiente}[1]{\textcolor{grispizarra}{\textit{[#1]}}}
% Hueco reservado para una gráfica o figura que aún no existe
\newcommand{\huecofigura}[2][5.2cm]{%
  \tikz{\draw[grispizarra!45, dashed, line width=0.5pt, rounded corners=2pt, fill=azulfondo]
          (0,0) rectangle (0.92\linewidth,#1);
        \node[text=grispizarra, font=\small\itshape, align=center, text width=0.8\linewidth]
          at (0.46\linewidth,#1/2) {#2};}}

% =========================================================================
%        ENTORNOS FORMALES (enunciados numerados, estilo artículo)
% =========================================================================
\usepackage{amsthm}

\newtheoremstyle{cienplano}{10pt}{10pt}{\itshape}{}%
  {\bfseries\color{azulprof}}{.}{0.6em}%
  {\thmname{#1}\thmnumber{ #2}\thmnote{\normalfont\upshape\ (#3)}}
\newtheoremstyle{ciendef}{10pt}{10pt}{\normalfont}{}%
  {\bfseries\color{azulprof}}{.}{0.6em}%
  {\thmname{#1}\thmnumber{ #2}\thmnote{\normalfont\upshape\ (#3)}}
\newtheoremstyle{cienobs}{8pt}{8pt}{\normalfont}{}%
  {\itshape\color{azulprof}}{.}{0.6em}%
  {\thmname{#1}\thmnumber{ #2}\thmnote{\normalfont\upshape\ (#3)}}

\theoremstyle{cienplano}
\newtheorem{proposicion}{Proposición}[section]
\newtheorem{teorema}[proposicion]{Teorema}
\newtheorem{resultado}[proposicion]{Resultado}
\newtheorem{conclusion}[proposicion]{Conclusión}
\theoremstyle{ciendef}
\newtheorem{definicion}[proposicion]{Definición}
\newtheorem{hipotesis}[proposicion]{Hipótesis}
\newtheorem{calculo}[proposicion]{Cálculo}
\theoremstyle{cienobs}
\newtheorem{observacion}[proposicion]{Observación}
\newtheorem*{demostracion}{Demostración}

% --- Referencia formalizada --------------------------------------------------
% Caja fina con fondo tenue y título sobre el borde. Cuerpo en cursiva y, al
% pie, la fuente completa con autor, título y número de cita.
% Uso: \begin{referencia}[título]{clave-bib} ... \end{referencia}
\newcounter{referencia}
\renewcommand{\thereferencia}{\theproposicion}
\newtcolorbox{referenciabox}[2]{%
  enhanced, sharp corners,
  colback=azulfondo, colframe=azulclaro, boxrule=0.4pt,
  left=10pt, right=10pt, top=7pt, bottom=6pt,
  before skip=14pt, after skip=12pt,
  attach boxed title to top left={yshift*=-\tcboxedtitleheight/2, xshift=8pt},
  boxed title style={colback=white, colframe=white, sharp corners, boxrule=0pt,
                     left=3pt, right=3pt, top=1pt, bottom=1pt},
  fonttitle=\bfseries\color{azulprof}, coltitle=azulprof,
  fontupper=\itshape,
  title={Referencia~\thereferencia\ifx\relax#1\relax\else\ \normalfont(#1)\fi},
  after upper={\par\vspace{5pt}{\color{azulclaro}\hrule height 0.4pt}\vspace{4pt}%
    {\upshape\footnotesize\color{grispizarra}\raggedleft
     \citeauthor{#2}, \citetitle{#2}~\cite{#2}.\par}}}
\newenvironment{referencia}[2][]
  {\par\stepcounter{proposicion}\refstepcounter{referencia}\begin{referenciabox}{#1}{#2}}
  {\end{referenciabox}}
\crefname{referencia}{referencia}{referencias}
\Crefname{referencia}{Referencia}{Referencias}

% --- Método (caja azulada para los fundamentos del tratamiento de datos) -----
\newcounter{metodo}
\renewcommand{\themetodo}{\theproposicion}
\newtcolorbox{metodobox}[1]{%
  enhanced, breakable, sharp corners,
  colback=azulfondo, colframe=azulclaro, boxrule=0.4pt,
  left=10pt, right=10pt, top=7pt, bottom=5pt,
  before skip=14pt, after skip=12pt,
  attach boxed title to top left={yshift*=-\tcboxedtitleheight/2, xshift=8pt},
  boxed title style={colback=white, colframe=white, sharp corners, boxrule=0pt,
                     left=3pt, right=3pt, top=1pt, bottom=1pt},
  fonttitle=\bfseries\color{azulprof}, coltitle=azulprof,
  title={Método~\themetodo\ifx\relax#1\relax\else\ \normalfont(#1)\fi}}
\newenvironment{metodo}[1][]
  {\par\stepcounter{proposicion}\refstepcounter{metodo}\begin{metodobox}{#1}}
  {\end{metodobox}}
\crefname{metodo}{método}{métodos}
\Crefname{metodo}{Método}{Métodos}

% --- Resultado destacado ----------------------------------------------------
\newtcolorbox{destacado}[1][]{%
  enhanced, sharp corners, frame hidden, colback=azulfondo,
  borderline west={1.6pt}{0pt}{turquesa},
  left=10pt, right=8pt, top=6pt, bottom=6pt,
  before skip=10pt, after skip=12pt,
  before upper={\ifx\relax#1\relax\else{\scshape\small\color{azulacero}#1}\par\vspace{5pt}\fi}}

% --- Tarea pendiente ----------------------------------------------------------
% Caja discontinua para las partes que aún hay que redactar o calcular.
\newtcolorbox{pendientebox}[1][]{%
  enhanced, breakable, sharp corners, colback=white, colframe=grispizarra!50,
  boxrule=0.5pt, borderline={0.5pt}{0pt}{grispizarra!50, dashed}, frame hidden,
  left=10pt, right=10pt, top=6pt, bottom=6pt, before skip=10pt, after skip=10pt,
  fontupper=\small\color{grispizarra},
  before upper={{\scshape\color{azulacero}Pendiente\ifx\relax#1\relax\else\ \textperiodcentered\ #1\fi}\par\vspace{2pt}}}

\crefname{proposicion}{proposición}{proposiciones}
\Crefname{proposicion}{Proposición}{Proposiciones}
\crefname{teorema}{teorema}{teoremas}
\Crefname{teorema}{Teorema}{Teoremas}
\crefname{definicion}{definición}{definiciones}
\Crefname{definicion}{Definición}{Definiciones}
\crefname{resultado}{resultado}{resultados}
\Crefname{resultado}{Resultado}{Resultados}
\crefname{conclusion}{conclusión}{conclusiones}
\crefname{hipotesis}{hipótesis}{hipótesis}
\crefname{calculo}{cálculo}{cálculos}
\crefname{observacion}{observación}{observaciones}
\Crefname{observacion}{Observación}{Observaciones}
\crefname{section}{sección}{secciones}
\Crefname{section}{Sección}{Secciones}
\crefname{subsection}{sección}{secciones}
\Crefname{subsection}{Sección}{Secciones}
\crefname{subsubsection}{sección}{secciones}
\crefname{table}{tabla}{tablas}
\Crefname{table}{Tabla}{Tablas}
\crefname{figure}{figura}{figuras}
\Crefname{figure}{Figura}{Figuras}
\crefname{equation}{ecuación}{ecuaciones}
\Crefname{equation}{Ecuación}{Ecuaciones}
\crefformat{equation}{ecuación~(#2#1#3)}
\Crefformat{equation}{Ecuación~(#2#1#3)}
\crefrangeformat{equation}{ecuaciones~(#3#1#4)--(#5#2#6)}
\crefmultiformat{equation}{ecuaciones~(#2#1#3)}{ y~(#2#1#3)}{, (#2#1#3)}{ y~(#2#1#3)}

% --- Estilo de gráficas de datos ----------------------------------------------
\pgfplotsset{
  datos/.style={
    tick label style={font=\scriptsize}, label style={font=\small},
    axis lines=left, axis line style={grispizarra},
    grid=major, grid style={grisclaro, line width=0.4pt},
    every tick/.style={grispizarra},
    legend style={font=\scriptsize, draw=none, fill=white, fill opacity=0.9,
                  text opacity=1, cells={anchor=west}},
    legend cell align=left,
    /pgf/number format/use period, /pgf/number format/1000 sep={\,},
    scaled ticks=false,
  },
  puntoA/.style={only marks, mark=*, mark size=2.1pt, azulprof,
    mark options={fill=azulprof, draw=white, line width=0.3pt}},
  puntoB/.style={only marks, mark=square*, mark size=1.9pt, azulacero,
    mark options={fill=azulclaro, draw=azulacero, line width=0.5pt}},
  puntoC/.style={only marks, mark=triangle*, mark size=2.3pt, turquesa!80!azulprof,
    mark options={fill=aguamarina, draw=azulpetroleo, line width=0.5pt}},
  teorica/.style={acento, line width=0.9pt, densely dashed},
  % familia de curvas de resonancia: del amortiguamiento más débil (marino)
  % al más fuerte (aguamarina)
  curva1/.style={azulprof, line width=0.8pt},
  curva2/.style={azulacero, line width=0.8pt},
  curva3/.style={azulpetroleo, line width=0.8pt},
  curva4/.style={turquesa, line width=0.8pt},
  curva5/.style={aguamarina!90!azulpetroleo, line width=0.8pt},
}

% Referencia a un panel de figura: «figura 2\,(a)», enlazada en bloque
\newcommand{\figpanel}[2]{\hyperref[#1]{figura~\ref*{#1}\,\textup{(#2)}}}

% Pies de figura y tabla
\usepackage[font=small, labelfont={sc,color=azulprof}, labelsep=period,
            margin=0.6cm, justification=justified]{caption}
\usepackage{subcaption}
\captionsetup[sub]{font=small, labelfont={sc,color=azulprof}, labelformat=parens,
                   margin=0pt, justification=centering}

% Fotografías con un marco fino
\newcommand{\foto}[2][\linewidth]{%
  \tcbox[enhanced, sharp corners, colback=white, colframe=grisclaro, boxrule=0.5pt,
         boxsep=0pt, left=0pt, right=0pt, top=0pt, bottom=0pt, nobeforeafter, valign=center]{%
    \includegraphics[width=\dimexpr#1-1pt\relax]{#2}}}

% =========================================================================
%                     VARIABLES PARA LA PORTADA
% =========================================================================
\newcommand{\Universidad}{Universidad de Salamanca}
\newcommand{\Facultad}{Facultad de Ciencias}
\newcommand{\Titulacion}{Grado en Física}
\newcommand{\Asignatura}{Laboratorio de Mecánica y Ondas}
\newcommand{\NumPractica}{Práctica 4}
\newcommand{\TituloPractica}{Péndulo de Pohl}
\newcommand{\SubtituloPractica}{Oscilaciones libres, amortiguadas y forzadas: la resonancia}
\newcommand{\tipotrabajo}{Informe de laboratorio}
\newcommand{\miNombre}{David Gómez García \\ Daniel Martín Guerrero}
\newcommand{\miTutor}{Mario Amado \\ Ivan de Martino \\ Manuel Pino}
\newcommand{\miUbicacion}{Salamanca}

% =========================================================================
%                     DATOS DEL ARTÍCULO (cabecera)
% =========================================================================
\title{Péndulo de Pohl}
\author{David Gómez García}
\author{Daniel Martín Guerrero}
\address{Grado en Física, Facultad de Ciencias, Universidad de Salamanca,
         Salamanca, España}
\date{}

% =========================================================================
%                                DOCUMENTO
% =========================================================================
\begin{document}
% --- PORTADA ----------------------------------------------------------------
\pagenumbering{alph}
%%%%%%%%%%%%%%%%%%%%%%%%%%%%%%%%%%%%%%%%%%%%%%%%%%%%%%%%%%%%%%%%%%%%%%%%
% Portada — estilo artículo científico, motivo: el péndulo de Pohl
% Misma estructura que las de los informes anteriores (cabecera con
% logotipos, motivo gráfico, franja, bloque de título y pie con autores),
% en una gama de azules que va del marino al aguamarina. El motivo es el
% aparato del laboratorio dibujado de frente; debajo, las dos mitades de
% la práctica: la oscilación libre que se apaga y las curvas de resonancia.
% Requiere (definidos en main.tex): colores azulprof, azulacero,
% azulpetroleo, turquesa, aguamarina, azulclaro, grispizarra, grisclaro,
% azulfondo; macros \Universidad, \Facultad, \Titulacion, \Asignatura,
% \NumPractica, \TituloPractica, \SubtituloPractica, \tipotrabajo,
% \miNombre, \miTutor, \miUbicacion; y figuras/pohl.tex.
%%%%%%%%%%%%%%%%%%%%%%%%%%%%%%%%%%%%%%%%%%%%%%%%%%%%%%%%%%%%%%%%%%%%%%%%

\begin{titlepage}
\thispagestyle{empty}
\pagecolor{white}

\begin{tikzpicture}[remember picture, overlay]
  \coordinate (O) at (current page.south west);

  % ----------------------------------------------------------------
  % Cabecera: logotipos y filete
  % ----------------------------------------------------------------
  \node[anchor=north west, inner sep=0] at ($(O)+(2.2cm,27.9cm)$)
    {\includegraphics[height=1.42cm]{figuras/logo-usal.png}};
  \node[anchor=north east, inner sep=0] at ($(O)+(18.8cm,27.9cm)$)
    {\includegraphics[height=1.42cm]{figuras/logo-ciencias.png}};
  \draw[grispizarra!60, line width=0.4pt]
    ($(O)+(2.2cm,26.0cm)$) -- ($(O)+(18.8cm,26.0cm)$);
  \node[anchor=north west, inner sep=0, text=grispizarra, font=\scshape\small]
    at ($(O)+(2.2cm,25.75cm)$)
    {\textls[80]{\Universidad\ \textbar\ \Facultad}};
  \node[anchor=north east, inner sep=0, text=grispizarra, font=\scshape\small]
    at ($(O)+(18.8cm,25.75cm)$)
    {\textls[80]{\Titulacion}};

  % ----------------------------------------------------------------
  % Motivo gráfico: el péndulo de Pohl de frente, con el anillo desviado
  % y el motor a media vuelta. Rótulos discretos, por fuera del aparato.
  % ----------------------------------------------------------------
  \begin{scope}[shift={($(O)+(9.55cm,20.05cm)$)}, x=0.88cm, y=0.88cm]
    % halo tenue detrás de la escala
    \fill[azulfondo] (0,0) circle (4.05);
    \draw[azulclaro, line width=0.4pt] (0,0) circle (4.05);
    \pohl{-28}{205}{0}
    % ángulo de giro
    \draw[grispizarra!60, densely dashed, line width=0.4pt] (0,0) -- (90:4.3);
    \draw[azulprof, line width=0.5pt, -{Stealth[length=4pt]}]
      (90:3.75) arc[start angle=90, end angle=62, radius=3.75];
    \node[text=azulprof] at (76:4.12) {$\varphi$};
    % rótulos
    \draw[grispizarra!70, line width=0.35pt] (150:0.85) -- (150:2.0) -- ++(-3.05,0);
    \fill[azulprof] (150:0.85) circle (1.2pt);
    \node[anchor=east, text=grispizarra, font=\itshape\small]
      at ($(150:2.0)+(-3.1,0)$) {muelle en espiral, $k$};
    \draw[grispizarra!70, line width=0.35pt] (200:2.22) -- (200:4.35) -- ++(-0.6,0);
    \fill[azulprof] (200:2.22) circle (1.2pt);
    \node[anchor=east, text=grispizarra, font=\itshape\small]
      at ($(200:4.35)+(-0.65,0)$) {anillo de cobre, $I$};
    \draw[grispizarra!70, line width=0.35pt] (-8:3.15) -- (-8:4.45) -- ++(0.6,0);
    \fill[white] (-8:3.15) circle (1.3pt);
    \node[anchor=west, text=grispizarra, font=\itshape\small]
      at ($(-8:4.45)+(0.65,0)$) {escala graduada};
    \draw[grispizarra!70, line width=0.35pt] (0.64,-2.25) -- (1.9,-2.25) -- (2.6,-1.55) -- ++(2.3,0);
    \fill[azulprof] (0.64,-2.25) circle (1.2pt);
    \node[anchor=west, text=grispizarra, font=\itshape\small]
      at (4.95,-1.55) {electroimán, $\gamma$};
    \draw[grispizarra!70, line width=0.35pt] ($(Motor)+(0,-0.62)$) -- ++(0,-0.38) -- ++(-1.05,0);
    \fill[azulprof] ($(Motor)+(0,-0.62)$) circle (1.2pt);
    \node[anchor=east, text=grispizarra, font=\itshape\small]
      at ($(Motor)+(-1.1,-1.0)$) {motor y biela, $\Omega$};
  \end{scope}

  % ----------------------------------------------------------------
  % Franja: las dos mitades de la práctica. A la izquierda, la oscilación
  % libre que se apaga dentro de su envolvente exponencial, con los puntos
  % de retorno; a la derecha, las curvas de resonancia para tres
  % amortiguamientos.
  % ----------------------------------------------------------------
  \begin{scope}[shift={($(O)+(0,13.35cm)$)}, x=1cm, y=1cm]
    % oscilación libre: phi0 e^{-g t} [cos(w t) + (g/w) sin(w t)]
    \def\xa{2.2}\def\xb{9.9}\def\gm{0.27}\def\wv{5.3}\def\Am{0.98}
    \draw[grispizarra!35, line width=0.3pt] (\xa,0) -- (\xb,0);
    \draw[azulacero!75, line width=0.5pt, densely dashed, domain=\xa:\xb, samples=80]
      plot (\x,{\Am*exp(-\gm*(\x-\xa))});
    \draw[azulacero!75, line width=0.5pt, densely dashed, domain=\xa:\xb, samples=80]
      plot (\x,{-\Am*exp(-\gm*(\x-\xa))});
    \draw[azulprof, line width=0.6pt, domain=\xa:\xb, samples=700, smooth]
      plot (\x,{\Am*exp(-\gm*(\x-\xa))*(cos(\wv*(\x-\xa) r)+(\gm/\wv)*sin(\wv*(\x-\xa) r))});
    \foreach \n in {0,...,13}{
      \pgfmathsetmacro{\tn}{\n*pi/\wv}
      \pgfmathsetmacro{\yn}{(-1)^\n*\Am*exp(-\gm*\tn)}
      \fill[turquesa, draw=white, line width=0.3pt] (\xa+\tn,\yn) circle (1.5pt);}
    \node[anchor=east, text=azulprof, font=\footnotesize] at (\xa-0.08,0) {$\varphi$};
    \node[anchor=south west, text=azulacero, font=\footnotesize, fill=white, inner sep=1.2pt]
      at (\xa+3.0,0.5) {$\pm\varphi_0e^{-\gamma t}$};
    \node[anchor=north, text=grispizarra, font=\scshape\scriptsize] at ({(\xa+\xb)/2},-1.32)
      {\textls[160]{la oscilación libre se apaga}};
    % separador
    \fill[azulclaro] (10.5,0) circle (1.4pt);
    % curvas de resonancia: u = Omega/omega0 en [0.35, 1.65]
    \def\xc{11.1}\def\xd{18.8}
    \pgfmathsetmacro{\escx}{(\xd-\xc)/1.3}
    \draw[grispizarra!35, line width=0.3pt] (\xc,-1.0) -- (\xd,-1.0);
    \draw[grispizarra!50, densely dotted, line width=0.4pt]
      ({\xc+0.65*\escx},-1.0) -- ({\xc+0.65*\escx},1.25);
    \foreach \g/\c in {0.1/azulprof, 0.18/azulacero, 0.32/turquesa}{
      \draw[\c, line width=0.6pt, domain=0.35:1.65, samples=220, smooth]
        plot ({\xc+(\x-0.35)*\escx},{-1.0+0.4*1/sqrt((1-\x*\x)^2+4*\g*\g*\x*\x)});}
    \node[anchor=north, text=grispizarra, font=\footnotesize] at ({\xc+0.65*\escx},-1.02) {$\wz$};
    \node[anchor=west, text=azulprof, font=\footnotesize] at ({\xc+0.73*\escx},0.95) {$\beta(\Omega)$};
    \node[anchor=north, text=grispizarra, font=\scshape\scriptsize] at ({(\xc+\xd)/2},-1.32)
      {\textls[160]{la forzada resuena}};
  \end{scope}

  % ----------------------------------------------------------------
  % Bloque del título
  % ----------------------------------------------------------------
  \node[anchor=north west, inner sep=0, text width=16.6cm, align=left]
    at ($(O)+(2.2cm,11.0cm)$)
    {{\color{azulacero}\scshape\large\textls[120]{\Asignatura\ \textbar\ \NumPractica}}\\[12pt]
     {\color{azulprof}\fontsize{36}{41}\selectfont\bfseries\spaceskip=0.28em\relax\TituloPractica}\\[10pt]
     {\color{azulpetroleo}\itshape\Large\SubtituloPractica}\\[12pt]
     {\color{turquesa}\rule{2.5cm}{1.4pt}}\\[9pt]
     {\color{grispizarra}\large
      $\ddot\varphi+2\gamma\dot\varphi+\wz^{2}\varphi=\alpha_0\sin\Omega t$
      \quad\textperiodcentered\quad
      $\WR=\sqrt{\wz^{2}-2\gamma^{2}}$}};

  % ----------------------------------------------------------------
  % Pie: tipo de trabajo, autores, tutores, lugar y fecha
  % ----------------------------------------------------------------
  \draw[grispizarra!60, line width=0.4pt]
    ($(O)+(2.2cm,4.6cm)$) -- ($(O)+(18.8cm,4.6cm)$);

  \node[anchor=north west, inner sep=0, text width=4.6cm, align=left]
    at ($(O)+(2.2cm,4.25cm)$)
    {{\scshape\small\color{azulacero}\textls[80]{Documento}}\\[3pt]
     {\large\bfseries\color{azulprof}\tipotrabajo}};

  \node[anchor=north west, inner sep=0, text width=5.6cm, align=left]
    at ($(O)+(7.2cm,4.25cm)$)
    {{\scshape\small\color{azulacero}\textls[80]{Autores}}\\[3pt]
     {\large\color{black!85}\miNombre}};

  \node[anchor=north west, inner sep=0, text width=5cm, align=left]
    at ($(O)+(13.8cm,4.25cm)$)
    {{\scshape\small\color{azulacero}\textls[80]{Tutores}}\\[3pt]
     {\large\color{black!85}\miTutor}};

  \node[anchor=south east, inner sep=0, text=grispizarra, font=\scshape\small]
    at ($(O)+(18.8cm,1.2cm)$)
    {\textls[100]{\miUbicacion, \number\year}};
  % franja inferior: la gama completa, del marino al aguamarina
  \fill[azulprof] ($(O)+(0,0)$) rectangle ($(O)+(21cm,0.45cm)$);
  \fill[azulacero] ($(O)+(0,0)$) rectangle ($(O)+(11cm,0.45cm)$);
  \fill[azulpetroleo] ($(O)+(0,0)$) rectangle ($(O)+(7cm,0.45cm)$);
  \fill[turquesa] ($(O)+(0,0)$) rectangle ($(O)+(4.2cm,0.45cm)$);
  \fill[aguamarina] ($(O)+(0,0)$) rectangle ($(O)+(2cm,0.45cm)$);
\end{tikzpicture}

\end{titlepage}

\clearpage
\pagenumbering{arabic}
\setcounter{page}{1}

\begin{abstract}
Estudiamos el péndulo de Pohl, un oscilador de torsión formado por un anillo
de cobre y un muelle en espiral, frenado por un electroimán y forzado por un
motor a través de una biela. En la oscilación libre, sin electroimán y con el
puente de la fuente en las posiciones~4 y~6, el periodo nos da la frecuencia
natural, $\wz=(\pend{$\cdots\pm\cdots$})\ \si{\radian\per\second}$, y el
decrecimiento exponencial de los máximos, linealizado como
$\ln\lvert\varphi_n\rvert$ frente a $t_n$, los coeficientes de
amortiguamiento $\gamma_0$, $\gamma_4$ y $\gamma_6$ (\pend{valores}). En la
oscilación forzada medimos la amplitud estacionaria $\beta$ para siete
frecuencias $\Omega$ en torno a la resonancia y los tres amortiguamientos, y
comparamos las curvas $\beta(\Omega)$ con la frecuencia de resonancia
esperada, $\WR=\sqrt{\wz^{2}-2\gamma^{2}}$. Los máximos se extraen con el
análisis de picos de PHYWE Measure y los resultados se dan con incertidumbre
típica y ampliada al \SI{95}{\percent}. \pend{Conclusión principal.}
\end{abstract}

\maketitle
\tableofcontents

% =========================================================================
\newpage
\seccion{Introducción}{Un oscilador que se deja frenar y empujar}\label{sec:introduccion}
% =========================================================================
Casi todo lo que oscila en la naturaleza lo hace en presencia de dos
influencias externas: algo que le roba energía y algo que, de vez en cuando,
se la devuelve. Un columpio se detiene por el rozamiento y vuelve a subir
cuando alguien lo empuja en el momento justo; la membrana de un altavoz, un
circuito sintonizado o la estructura de un puente responden de forma muy
distinta según la frecuencia con que se les excite. Entender ese equilibrio
entre disipación y forzamiento, y en particular el fenómeno de la
\emph{resonancia}, es uno de los resultados centrales de la física de las
oscilaciones~\cite{FIV,FRENCH}.

El péndulo que estudiamos lo ideó Robert Wichard Pohl en Gotinga, a comienzos
del siglo~XX, precisamente para mostrar ese fenómeno en el aula, y su diseño
apenas ha cambiado desde entonces. Es un oscilador de torsión: un anillo de
cobre que gira alrededor de su eje, recuperado por un muelle en espiral. Su
interés está en que las dos influencias externas se pueden controlar por
separado y sin tocar el oscilador. Un electroimán induce en el anillo
corrientes de Foucault que lo frenan con un momento proporcional a su
velocidad angular, y cuya intensidad fijamos con la tensión de la bobina; un
motor, unido al muelle por una biela, lo empuja periódicamente con la
frecuencia que elijamos~\cite{GUION-P4}. Así podemos observar primero cómo se
apaga una oscilación libre y, después, cómo responde el anillo cuando se le
fuerza cerca de su frecuencia natural.

\subsection{Objetivos}\label{sec:objetivos}
\paragraph{Oscilación libre amortiguada.}
\begin{itemize}
  \item Medir el periodo de la oscilación libre y deducir de él la frecuencia
        natural $\wz$ del péndulo, con su incertidumbre.
  \item Comprobar que las amplitudes máximas decrecen exponencialmente y
        determinar el coeficiente de amortiguamiento $\gamma$ sin electroimán
        y con el puente de la fuente en las posiciones~4 y~6, mediante un
        ajuste lineal de $\ln\lvert\varphi_n\rvert$ frente a $t_n$.
\end{itemize}
\paragraph{Oscilación forzada.}
\begin{itemize}
  \item Medir la amplitud estacionaria $\beta$ para varias frecuencias de
        forzamiento $\Omega$ en torno a la resonancia y para los tres grados
        de amortiguamiento, y representar las curvas de resonancia.
  \item Calcular, a partir de $\wz$ y $\gamma$, la frecuencia de resonancia
        esperada $\WR$ para cada amortiguamiento, con su incertidumbre, y
        contrastarla con las curvas medidas.
  \item Discutir cómo cambian la altura, la anchura y la posición del pico de
        resonancia con el amortiguamiento.
\end{itemize}
Para ello nos apoyaremos en las \cref{eq:linealizacion,eq:beta,eq:resonancia}
de la \cref{sec:fundamentos}.

\subsection{Material empleado}\label{sec:material}
Hemos utilizado el siguiente instrumental (\cref{fig:montaje,fig:auxiliar}):
\begin{itemize}
  \item Péndulo de Pohl: anillo de cobre que gira en torno a un eje horizontal,
        unido a un \emph{muelle en espiral}; fiel solidario con el anillo y
        escala graduada fija para leer la desviación angular.
  \item Electroimán de frenado, cuyas bobinas abrazan la parte inferior del
        anillo.
  \item Motor de corriente continua con disco excéntrico y biela, y dos
        potenciómetros para ajustar su velocidad: uno de ajuste grueso y otro
        de ajuste fino.
  \item Fuente de alimentación universal, con un puente enchufable que
        selecciona la tensión de salida, y rectificador de corriente.
  \item Unidad de medida Cobra-3 y ordenador con el programa PHYWE
        Measure~\cite{PHYWE}.
  \item Cables de conexión.
\end{itemize}

\subsection{Montaje experimental}\label{sec:montaje}
La \cref{fig:montaje} muestra el aparato en el laboratorio. El anillo de cobre
está unido al eje de giro por cuatro radios y, a través del eje, al extremo
interior del muelle en espiral; el extremo exterior del muelle está sujeto a
una palanca, el \emph{excitador}, que en la oscilación libre permanece
inmóvil. Cuando el motor gira, el disco excéntrico arrastra la biela, la biela
hace oscilar la palanca y esta, a su vez, retuerce periódicamente el muelle:
así es como el motor le comunica al anillo un momento de fuerzas periódico.
El fiel, solidario con el anillo, marca sobre la escala la desviación $\varphi$
respecto de la posición de equilibrio, que debe coincidir con el cero; si no
lo hace, se corrige girando a mano el disco del motor.

El electroimán se alimenta desde la fuente a través del rectificador, que
convierte en continua la tensión alterna seleccionada con el puente enchufable
(\cref{fig:auxiliar}). Sin el puente no circula corriente por la bobina y solo
queda el rozamiento propio del aparato; cuanto mayor es la tensión elegida,
más intenso es el campo y más fuerte el frenado. Un sensor en el eje del
péndulo convierte el ángulo de giro en una señal que la unidad Cobra-3
digitaliza y que PHYWE Measure representa en tiempo real como $\varphi(t)$.
La \cref{fig:esquema} resume el funcionamiento del aparato y los momentos que
actúan sobre el anillo.

% --- Fotografía rotulada del montaje ---------------------------------------------
\begin{figure}[!htbp]
\centering
\begin{tikzpicture}[x=1cm, y=1cm, >={Stealth[length=4pt]},
    rotulo/.style={font=\small, text=azulprof, align=left, inner sep=1.5pt},
    rotuloD/.style={font=\small, text=azulprof, align=right, inner sep=1.5pt},
    guia/.style={azulacero, line width=0.5pt},
    blanco/.style={circle, fill=white, draw=azulprof, line width=0.6pt, inner sep=0pt, minimum size=4.6pt}]
  % la foto ocupa 9 cm de ancho (2100 x 2268 px) y su esquina inferior izquierda está en (0,0)
  \def\W{9}\def\H{9.72}
  \node[anchor=south west, inner sep=0] (F) at (0,0) {\foto[9cm]{figuras/montaje-pohl.jpg}};
  % coordenadas relativas (x/W, y/H) medidas sobre la foto, con y desde arriba
  \newcommand{\punto}[3]{\coordinate (#1) at ({#2*\W},{(1-#3)*\H});}
  \punto{pFiel}{0.228}{0.158}
  \punto{pEscala}{0.072}{0.172}
  \punto{pAnillo}{0.133}{0.212}
  \punto{pMuelle}{0.205}{0.305}
  \punto{pIman}{0.265}{0.440}
  \punto{pBiela}{0.300}{0.621}
  \punto{pFuente}{0.900}{0.370}
  \punto{pExcen}{0.500}{0.690}
  \punto{pRect}{0.950}{0.731}
  \punto{pPot}{0.745}{0.727}
  % columna izquierda de rótulos
  \foreach \p/\y/\t in {pFiel/9.30/{Fiel},
                        pEscala/8.30/{Escala graduada},
                        pAnillo/7.30/{Anillo de cobre},
                        pMuelle/6.30/{Muelle en espiral},
                        pIman/5.05/{Electroimán\\[-1pt]{\footnotesize(corrientes de Foucault)}},
                        pBiela/3.15/{Biela y palanca\\[-1pt]{\footnotesize del excitador}}}{
    \node[rotuloD, anchor=east] (r\p) at (-0.45,\y) {\t};
    \draw[white, line width=1.7pt, opacity=0.85] ($(r\p.east)+(0.25,0)$) -- (\p);
    \draw[guia] (r\p.east) -- ++(0.25,0) -- (\p);
    \node[blanco] at (\p) {};}
  % columna derecha de rótulos
  \foreach \p/\y/\t in {pFuente/7.30/{Fuente de\\[-1pt]alimentación},
                        pExcen/5.10/{Motor con disco\\[-1pt]excéntrico},
                        pRect/3.45/{Rectificador},
                        pPot/1.75/{Potenciómetros\\[-1pt]{\footnotesize(grueso y fino)}}}{
    \node[rotulo, anchor=west] (r\p) at (\W+0.45,\y) {\t};
    \draw[white, line width=1.7pt, opacity=0.85] ($(r\p.west)+(-0.25,0)$) -- (\p);
    \draw[guia] (r\p.west) -- ++(-0.25,0) -- (\p);
    \node[blanco] at (\p) {};}
\end{tikzpicture}
\caption{El péndulo de Pohl en el laboratorio. El anillo de cobre gira frente a
la escala graduada fija y su desviación se lee con el fiel. El muelle en espiral
une el eje del anillo con la palanca del excitador, que la biela mueve cuando
el motor gira. El electroimán, alimentado desde la fuente a través del
rectificador, frena el anillo induciendo en él corrientes de Foucault.}
\label{fig:montaje}
\end{figure}

% --- Material auxiliar -------------------------------------------------------------
\begin{figure}[!p]
\centering
\begin{subfigure}[b]{0.32\textwidth}\centering
  \foto{figuras/cobra3.jpg}
  \caption{Unidad Cobra-3.}
\end{subfigure}\hfill
\begin{subfigure}[b]{0.32\textwidth}\centering
  \foto{figuras/puente.jpg}
  \caption{Puente enchufable.}
\end{subfigure}\hfill
\begin{subfigure}[b]{0.32\textwidth}\centering
  \foto{figuras/fuente-rectificador.jpg}
  \caption{Fuente y rectificador.}
\end{subfigure}
\caption{Material auxiliar. (a)~La unidad Cobra-3 digitaliza la señal del
sensor de ángulo y la envía al ordenador. (b)~El puente enchufable de la fuente
selecciona la tensión de salida y, con ella, el grado de amortiguamiento.
(c)~El rectificador, un puente de diodos, convierte esa tensión en continua
antes de llevarla a la bobina del electroimán.}
\label{fig:auxiliar}

\vspace{1.4\baselineskip}

% --- Esquema del aparato y de los momentos que actúan sobre el anillo -----------
\begin{tikzpicture}[x=0.86cm, y=0.86cm, >={Stealth[length=5pt]}, font=\small]
  \pohl{24}{150}{1}
  % ángulo de giro: desde la vertical (equilibrio) hasta el fiel
  \draw[grispizarra!70, densely dashed, line width=0.4pt] (0,0) -- (90:3.95);
  \draw[azulprof, line width=0.5pt, ->] (90:3.75) arc[start angle=90, end angle=114, radius=3.75];
  \node[text=azulprof] at (102:4.15) {$\varphi$};
  % momento recuperador del muelle (sentido contrario al giro)
  \draw[acento, line width=0.9pt, ->] (62:1.3) arc[start angle=62, end angle=14, radius=1.3];
  \node[text=acento, fill=white, inner sep=1pt] at (36:1.62) {$-k\varphi$};
  % frenado magnético, junto al electroimán
  \draw[acento, line width=0.9pt, ->] (-42:1.72) arc[start angle=-42, end angle=-72, radius=1.72];
  \node[text=acento, fill=white, inner sep=1pt] at (-52:1.2) {$-\eta\dot\varphi$};
  \draw[grispizarra!70, line width=0.35pt] (0.64,-2.3) -- (3.7,-2.3);
  \fill[azulprof] (0.64,-2.3) circle (1.3pt);
  \node[text=grispizarra, font=\footnotesize\itshape, anchor=west] at (3.75,-2.3) {campo $\vec B$ del electroimán};
  % forzamiento
  \draw[azulacero, line width=0.5pt, ->] ($(Motor)+(0.85,0.3)$) arc[start angle=20, end angle=120, radius=0.9];
  \node[text=azulacero, anchor=south] at ($(Motor)+(0.35,0.95)$) {$\Omega$};
  \node[text=azulprof, anchor=north east, align=right] at ($(Palanca)+(-0.1,-0.2)$)
    {$M(t)\simeq M_0\sin\Omega t$};
  % rótulos discretos
  \draw[grispizarra!70, line width=0.35pt] (135:2.22) -- (135:4.25) -- ++(-0.5,0);
  \node[text=grispizarra, font=\footnotesize\itshape, anchor=east] at ($(135:4.25)+(-0.55,0)$) {anillo de cobre, $I$};
  \draw[grispizarra!70, line width=0.35pt] (160:0.62) -- (160:1.55) -- ++(-3.1,0);
  \node[text=grispizarra, font=\footnotesize\itshape, anchor=east] at ($(160:1.55)+(-3.15,0)$) {muelle en espiral, $k$};
  \draw[grispizarra!70, line width=0.35pt] (35:3.2) -- (35:4.1) -- ++(0.5,0);
  \node[text=grispizarra, font=\footnotesize\itshape, anchor=west] at ($(35:4.1)+(0.55,0)$) {escala fija};
\end{tikzpicture}
\caption{Esquema del péndulo de Pohl, de frente y sin escala. El anillo, de
momento de inercia $I$, se ha desviado un ángulo $\varphi$ desde la vertical
(discontinua). Sobre él actúan el momento recuperador del muelle, $-k\varphi$,
y el del frenado magnético, $-\eta\dot\varphi$, debido a las corrientes que el
campo $\vec B$ del electroimán induce en el cobre (cruces: el campo entra en el
papel). Con el motor en marcha, la biela mueve la palanca del excitador y el
muelle transmite al anillo un momento aproximadamente armónico, $M(t)$.}
\label{fig:esquema}
\end{figure}

\ornamento

% =========================================================================
\newpage
\seccion{Fundamentos teóricos}{Disipación, forzamiento y resonancia}\label{sec:fundamentos}
% =========================================================================
\subsection{Ecuación del movimiento}\label{sec:ecuacion}
El anillo es un sólido rígido que gira alrededor de un eje fijo, de modo que su
dinámica se reduce a una sola ecuación: el momento de inercia $I$ por la
aceleración angular $\ddot\varphi$ es igual a la suma de los momentos de las
fuerzas respecto del eje. Sobre él actúan tres:
\begin{itemize}
  \item \emph{El muelle en espiral.} Para deformaciones moderadas cumple la
        ley de Hooke en su versión de torsión: ejerce un momento
        $-k\varphi$, proporcional al ángulo girado y de sentido contrario, con
        $k$ la constante de torsión del muelle.
  \item \emph{El frenado magnético.} Al moverse la parte inferior del anillo
        dentro del campo $\vec B$ del electroimán, en el cobre se induce una
        fuerza electromotriz y circulan corrientes de Foucault. Por la ley de
        Lenz, la fuerza que el campo ejerce sobre esas corrientes se opone al
        movimiento que las crea y, como la densidad de corriente es
        proporcional a la velocidad, el momento resultante es
        $-\eta\dot\varphi$~\cite{PURCELL}. El coeficiente $\eta$ crece con el
        cuadrado del campo y, por tanto, con la corriente de la bobina: esa es
        la palanca con la que regulamos el amortiguamiento. A él se suma un
        pequeño rozamiento mecánico, que supondremos también proporcional a la
        velocidad.
  \item \emph{El forzamiento.} Con el motor en marcha, la biela hace oscilar el
        extremo exterior del muelle, que transmite al anillo un momento
        periódico de la misma frecuencia angular $\Omega$ que el motor. Como
        la biela es mucho más larga que el radio de la excéntrica, ese momento
        es prácticamente armónico, $M(t)=M_0\sin\Omega t$~\cite{GUION-P4}.
\end{itemize}
La segunda ley de Newton para la rotación da, por tanto,
\begin{equation}\label{eq:newton}
  I\ddot\varphi=-k\varphi-\eta\dot\varphi+M_0\sin\Omega t .
\end{equation}

\begin{definicion}[Parámetros del oscilador]\label{def:parametros}
Dividiendo la \cref{eq:newton} por $I$ y definiendo
\begin{equation}\label{eq:parametros}
  \wz\equiv\sqrt{\frac{k}{I}},\qquad
  \gamma\equiv\frac{\eta}{2I},\qquad
  \alpha_0\equiv\frac{M_0}{I},
\end{equation}
la ecuación del movimiento adopta la forma canónica del oscilador armónico
amortiguado y forzado:
\begin{equation}\label{eq:movimiento}
  \ddot\varphi+2\gamma\dot\varphi+\wz^{2}\varphi=\alpha_0\sin\Omega t .
\end{equation}
$\wz$ es la \emph{frecuencia natural}, con la que oscilaría el anillo sin
rozamiento; $\gamma$, el \emph{coeficiente de amortiguamiento}, con
dimensiones de inversa de tiempo; y $\alpha_0$ mide la intensidad del
forzamiento\footnote{En los apuntes de Física~IV~\cite{FIV} el término
disipativo se escribe $\gamma\dot x$, de modo que aquel coeficiente es el doble
del nuestro. Aquí seguimos la convención del guion~\cite{GUION-P4}, con el
factor~2 explícito, que simplifica las soluciones.}.
\end{definicion}

\subsection{Oscilación libre amortiguada}\label{sec:libre}
Con el motor parado ($\alpha_0=0$), la \cref{eq:movimiento} es lineal y
homogénea, con coeficientes constantes. Ensayando soluciones
$\varphi\propto e^{\lambda t}$ se llega a la ecuación característica
$\lambda^{2}+2\gamma\lambda+\wz^{2}=0$, cuyas raíces son
\begin{equation}\label{eq:raices}
  \lambda_{\pm}=-\gamma\pm\sqrt{\gamma^{2}-\wz^{2}} .
\end{equation}
El signo del discriminante distingue tres regímenes~\cite{FIV}: amortiguamiento
débil o subcrítico ($\gamma<\wz$), crítico ($\gamma=\wz$) y supercrítico
($\gamma>\wz$). Solo en el primero el anillo llega a oscilar, y es el único que
exploramos en esta práctica.

\begin{proposicion}[Oscilación subamortiguada]\label{prop:libre}
Si $\gamma<\wz$, las raíces son complejas conjugadas y la solución real general
es
\begin{equation}\label{eq:libre}
  \varphi(t)=A_0\,e^{-\gamma t}\sin(\wa t+\phi_0),\qquad
  \wa=\sqrt{\wz^{2}-\gamma^{2}},
\end{equation}
donde $A_0$ y $\phi_0$ dependen de las condiciones iniciales. El anillo oscila
con una frecuencia $\wa$ algo menor que la natural, y su amplitud se apaga
exponencialmente con un tiempo característico $1/\gamma$.
\end{proposicion}

Las magnitudes que medimos directamente son los \emph{puntos de retorno}, en
los que el anillo se detiene un instante antes de invertir su giro. Si lo
soltamos desde el reposo con una desviación $\varphi_0$, imponer
$\dot\varphi(0)=0$ fija $\tan\phi_0=\wa/\gamma$, y la velocidad se anula cada
medio periodo, en los instantes $t_n=n\pi/\wa$. En ellos
\begin{equation}\label{eq:retornos}
  \varphi_n\equiv\varphi(t_n)=(-1)^{n}\varphi_0\,e^{-\gamma t_n},\qquad n=0,1,2,\dots
\end{equation}
Los máximos de un mismo lado, que son los que localiza el programa, están
separados por el \emph{pseudoperiodo} $T=2\pi/\wa$.

\begin{proposicion}[Linealización del decrecimiento]\label{prop:linealizacion}
Tomando logaritmos en el valor absoluto de la \cref{eq:retornos},
\begin{equation}\label{eq:linealizacion}
  \ln\lvert\varphi_n\rvert=\ln\lvert\varphi_0\rvert-\gamma\,t_n
  \qquad\Longleftrightarrow\qquad
  \ln\left\lvert\frac{\varphi_0}{\varphi_n}\right\rvert=\gamma\,t_n .
\end{equation}
Los puntos $(t_n,\ln\lvert\varphi_n\rvert)$ caen, por tanto, sobre una recta
de pendiente $-\gamma$: un ajuste lineal da directamente el coeficiente de
amortiguamiento, y la propia linealidad de los datos pone a prueba la hipótesis
de un frenado proporcional a la velocidad.
\end{proposicion}

\begin{observacion}[Lo que revela la forma del decrecimiento]\label{obs:coulomb}
Si el frenado dominante fuera un rozamiento seco, de módulo constante, las
amplitudes decrecerían en progresión \emph{aritmética} y no geométrica: una
cantidad fija por ciclo~\cite{TAYLOR-CM}. En la representación
semilogarítmica de la \cref{eq:linealizacion}, eso se manifestaría como una
curvatura hacia abajo al final del registro, cuando las amplitudes son
pequeñas. Vigilar esa curvatura, sobre todo sin electroimán, nos dirá hasta
dónde es válido el modelo.
\end{observacion}

\begin{observacion}[Frecuencia natural y frecuencia medida]\label{obs:w0}
El periodo que medimos es el de la oscilación amortiguada, de modo que nos da
$\wa$ y no $\wz$. Ambas se relacionan por la \cref{eq:libre},
$\wz=\sqrt{\wa^{2}+\gamma^{2}}$, y su diferencia relativa,
$(\wz-\wa)/\wz\simeq\gamma^{2}/(2\wz^{2})$, es de segundo orden en
$\gamma/\wz$. Bastará, pues, con medir el periodo sin electroimán y corregirlo
con el $\gamma_0$ obtenido, comprobando que la corrección es despreciable
frente a la incertidumbre.
\end{observacion}

Dos cantidades adimensionales resumen la intensidad del amortiguamiento. El
\emph{decremento logarítmico}, $\Lambda=\gamma T$, es el logaritmo del cociente
entre dos máximos consecutivos de un mismo lado. El \emph{factor de calidad},
que con nuestra convención es $Q=\wz/(2\gamma)$, indica cuántos radianes de
oscilación tarda la energía, que decae como $e^{-2\gamma t}$, en reducirse en
un factor $e$; para amortiguamiento débil, la energía perdida en cada ciclo es
una fracción $2\pi/Q$ de la almacenada~\cite{FIV}.

\subsection{Oscilación forzada: transitorio y estado estacionario}\label{sec:forzada}
Con el motor en marcha, la solución general de la \cref{eq:movimiento} es la
suma de la solución general de la ecuación homogénea, ya conocida, y de una
solución particular. Esta última se obtiene con comodidad pasando al plano
complejo: si $z(t)$ resuelve
$\ddot z+2\gamma\dot z+\wz^{2}z=\alpha_0e^{\ii\Omega t}$, su parte imaginaria
resuelve la \cref{eq:movimiento}. Ensayando $z=B\,e^{\ii\Omega t}$ resulta
$B=\alpha_0/(\wz^{2}-\Omega^{2}+2\ii\gamma\Omega)$, cuyo módulo y argumento dan
la amplitud y el desfase de la respuesta.

\begin{proposicion}[Respuesta a un forzamiento armónico]\label{prop:forzada}
Si $\gamma>0$, la solución de la \cref{eq:movimiento} es
\begin{equation}\label{eq:forzada}
  \varphi(t)=
  \underbrace{A_0\,e^{-\gamma t}\sin(\wa t+\phi_0)}_{\text{transitorio}}
  +\underbrace{\beta\sin(\Omega t-\delta)}_{\text{estacionario}},
\end{equation}
con
\begin{equation}\label{eq:beta}
  \beta(\Omega)=\frac{\alpha_0}{\sqrt{\left(\wz^{2}-\Omega^{2}\right)^{2}+4\gamma^{2}\Omega^{2}}},
  \qquad
  \tan\delta=\frac{2\gamma\Omega}{\wz^{2}-\Omega^{2}},\quad \delta\in[0,\pi].
\end{equation}
\end{proposicion}

El transitorio recuerda las condiciones iniciales y se apaga con el tiempo
característico $1/\gamma$; el término estacionario, en cambio, no depende de
ellas: el anillo acaba oscilando con la frecuencia del motor, con una amplitud
$\beta$ y un retraso $\delta$ que solo dependen de $\Omega$, $\wz$ y $\gamma$.
El retraso crece con la frecuencia: es casi nulo cuando el motor gira despacio,
vale exactamente $\pi/2$ cuando $\Omega=\wz$ (el anillo pasa por el
equilibrio cuando la biela está en un extremo) y tiende a $\pi$, oposición de
fase, cuando $\Omega\gg\wz$ (\figpanel{fig:resonancia}{b}).

\begin{observacion}[Cuánto hay que esperar]\label{obs:espera}
Para que el transitorio haya caído por debajo del \SI{1}{\percent} de su valor
inicial hace falta que $e^{-\gamma t}<10^{-2}$, es decir, $t>\ln 100/\gamma\approx
4.6/\gamma$. El tiempo de espera es, por tanto, más largo cuanto \emph{menor}
es el amortiguamiento: precisamente sin electroimán, donde la resonancia es más
aguda, es donde más paciencia hay que tener antes de empezar a registrar, y
donde una medida precipitada mezcla en la amplitud la contribución del
transitorio.
\end{observacion}

\subsection{La resonancia}\label{sec:resonancia}
Fijado el amortiguamiento, la amplitud estacionaria $\beta(\Omega)$ de la
\cref{eq:beta} es máxima cuando el radicando del denominador es mínimo.
Derivando respecto de $\Omega$ e igualando a cero,
$-4\Omega(\wz^{2}-\Omega^{2})+8\gamma^{2}\Omega=0$, lo que conduce al
resultado siguiente~\cite{GUION-P4,FIV}.

\begin{proposicion}[Frecuencia y amplitud de resonancia]\label{prop:resonancia}
Si $0<\gamma<\wz/\sqrt2$, la amplitud estacionaria crece con $\Omega$ desde
$\beta(0)=\alpha_0/\wz^{2}$ hasta un máximo, y luego decrece. El máximo se
alcanza en la \emph{frecuencia de resonancia}
\begin{equation}\label{eq:resonancia}
  \WR=\sqrt{\wz^{2}-2\gamma^{2}},
\end{equation}
y vale
\begin{equation}\label{eq:amplitudR}
  \bR=\frac{\alpha_0}{2\gamma\sqrt{\wz^{2}-\gamma^{2}}} .
\end{equation}
Si $\gamma\geq\wz/\sqrt2$, no hay pico: la amplitud decrece monótonamente desde
$\beta(0)$.
\end{proposicion}

La \cref{fig:resonancia} muestra la familia de curvas $\beta(\Omega)$ para
varios amortiguamientos. Tres rasgos suyos guiarán la discusión de nuestras
medidas:
\begin{itemize}
  \item \emph{La altura.} Para amortiguamiento débil,
        $\bR/\beta(0)\simeq\wz/(2\gamma)=Q$: el pico crece como $1/\gamma$ y,
        sin rozamiento, la amplitud de resonancia divergería.
  \item \emph{La anchura.} La curva de energía media, proporcional a
        $\beta^{2}$, tiene una anchura a media altura
        $\Delta\Omega\simeq2\gamma=\wz/Q$~\cite{FIV}: el pico se ensancha en
        proporción directa al amortiguamiento.
  \item \emph{La posición.} El desplazamiento del máximo respecto de la
        frecuencia natural, $\wz-\WR\simeq\gamma^{2}/\wz$, es de segundo orden
        en $\gamma/\wz$. Con amortiguamientos débiles, la resonancia delata el
        rozamiento por su altura y su anchura, mucho antes que por su
        posición.
\end{itemize}

\begin{figure}[!htbp]
\centering
\begin{tikzpicture}
\begin{groupplot}[
  group style={group size=2 by 1, horizontal sep=1.55cm},
  datos, width=0.5\textwidth, height=0.39\textwidth,
  xmin=0, xmax=2, domain=0.001:2, samples=260,
  xlabel={$\Omega/\wz$},
  xtick={0,0.5,1,1.5,2},
]
\nextgroupplot[ymin=0, ymax=7.2, ylabel={$\beta\,\wz^{2}/\alpha_0$},
  legend style={at={(0.02,0.98)}, anchor=north west},
  title={\small(a)\ Amplitud}, title style={text=azulprof, yshift=-4pt}]
  \foreach \g/\c in {0.08/curva1, 0.12/curva2, 0.2/curva3, 0.35/curva4, 0.707/curva5}{
    \edef\tmp{\noexpand\addplot[\c] {1/sqrt((1-x^2)^2+4*\g^2*x^2)};}\tmp}
  \addplot[grispizarra, densely dotted, line width=0.6pt, domain=0.001:0.9935]
    {1/sqrt(2*(1-x^4))};
  \legend{$\gamma/\wz=0.08$, $0.12$, $0.20$, $0.35$, $1/\sqrt2$, lugar de los máximos}
\nextgroupplot[ymin=0, ymax=180, ytick={0,45,90,135,180},
  ylabel={$\delta$ (\si{\degree})},
  title={\small(b)\ Desfase}, title style={text=azulprof, yshift=-4pt}]
  \foreach \g/\c in {0.08/curva1, 0.12/curva2, 0.2/curva3, 0.35/curva4, 0.707/curva5}{
    \edef\tmp{\noexpand\addplot[\c] {atan2(2*\g*x, 1-x^2)};}\tmp}
  \addplot[grispizarra!60, densely dashed, line width=0.4pt] coordinates {(0,90) (2,90)};
\end{groupplot}
\end{tikzpicture}
\caption{Respuesta estacionaria del oscilador forzado, según la
\cref{eq:beta}, para cinco amortiguamientos. (a)~Amplitud normalizada por su
valor estático $\alpha_0/\wz^{2}$: al crecer $\gamma$, el pico baja, se ensancha
y se desplaza, muy poco, hacia frecuencias menores. La línea punteada es el
lugar geométrico de los máximos, $\bR\wz^{2}/\alpha_0=[2(1-x^{4})]^{-1/2}$ con
$x=\WR/\wz$, que se obtiene eliminando $\gamma$ entre las
\cref{eq:resonancia,eq:amplitudR}; para $\gamma=\wz/\sqrt2$ el máximo se ha
desplazado hasta $\Omega=0$. (b)~Desfase entre el forzamiento y la respuesta:
todas las curvas pasan por $\delta=\pi/2$ en $\Omega=\wz$, y la transición es
tanto más brusca cuanto menor es el amortiguamiento.}
\label{fig:resonancia}
\end{figure}

\ornamento

% =========================================================================
\newpage
\seccion{Procedimiento experimental}{Cómo medimos}\label{sec:procedimiento}
% =========================================================================
Seguimos la metodología del guion~\cite{GUION-P4}, con las herramientas de
adquisición y análisis de PHYWE Measure~\cite{PHYWE} que describimos a
continuación.

\subsection{Preparación y registro}\label{sec:registro}
Con la fuente, la unidad Cobra-3 y el ordenador encendidos, abrimos PHYWE
Measure y elegimos el módulo de medida \emph{Traslación/Rotación}, que
registra el ángulo de giro del sensor frente al tiempo, comprobando que las
opciones de medida son las que indica el guion. Al pulsar \emph{Continuar}
aparece una ventana con la lectura instantánea del ángulo, que debe marcar
\SI{0}{\degree} con el anillo en reposo; si el fiel no apunta al cero de la
escala, se corrige girando el disco del motor. Cada registro se inicia con
\emph{Iniciar medida} y se detiene con \emph{Finalizar medida}, tras lo cual el
programa dibuja la curva $\varphi(t)$ (\cref{fig:registro}), que guardamos para
analizarla después.

\begin{figure}[!htbp]
\centering
\foto[0.82\textwidth]{figuras/registro-libre.jpg}
\caption{Registro de una oscilación libre amortiguada en PHYWE Measure: ángulo
de giro, en grados, frente al tiempo, en segundos. El anillo se suelta desde
el reposo con una desviación próxima a \SI{90}{\degree} y la amplitud se va
apagando ciclo a ciclo, mientras el periodo se mantiene constante, como
predice la \cref{eq:libre}.}
\label{fig:registro}
\end{figure}

\subsection{Oscilación libre amortiguada}\label{sec:proc-libre}
\begin{enumerate}
  \item \emph{Frecuencia natural.} Sin el puente enchufable, de modo que solo
        actúa el rozamiento propio del aparato, desplazamos el anillo a mano
        hasta una desviación grande, lo soltamos desde el reposo, sin darle
        impulso, e iniciamos el registro. Sin electroimán las oscilaciones
        duran mucho, lo que permite medir el periodo sobre un número elevado de
        ciclos.
  \item \emph{Amortiguamiento.} Repetimos el registro con el puente en las
        posiciones~4 y~6, sin cambiar nada más. En cada caso dejamos que el
        anillo oscile hasta que la amplitud sea pequeña.
\end{enumerate}
Los tres registros nos dan, para cada amortiguamiento, los instantes $t_n$ y
las amplitudes $\varphi_n$ de los máximos sucesivos, que son los datos de la
\cref{eq:linealizacion}.

\subsection{Oscilación forzada}\label{sec:proc-forzada}
\begin{enumerate}
  \item Retiramos el puente, encendemos el motor y situamos el potenciómetro
        de ajuste grueso entre~20 y~30, donde se encuentra la resonancia. Con
        los dos potenciómetros buscamos la frecuencia a la que la amplitud es
        máxima. Cerca de ella, sin electroimán, la amplitud llega a rebasar el
        recorrido de la escala, por lo que conviene acercarse con cuidado.
  \item Alcanzada la resonancia, esperamos a que se establezca el régimen
        estacionario (\cref{obs:espera}) y registramos más de veinte
        oscilaciones, de las que obtendremos la amplitud $\beta$ y la
        frecuencia $\Omega$. Sin tocar los potenciómetros, colocamos el puente
        en las posiciones~4 y~6 y, en cada caso, esperamos de nuevo al
        estacionario y registramos.
  \item Retiramos el puente y variamos ligeramente $\Omega$ con el
        potenciómetro fino. Repetimos el registro sin puente y con el puente
        en~4 y en~6.
  \item Repetimos el paso anterior hasta completar tres frecuencias por
        encima de la de resonancia y tres por debajo, modificando $\Omega$
        únicamente con el potenciómetro fino. En total son siete frecuencias
        y tres amortiguamientos, veintiún registros.
\end{enumerate}

\begin{pendientebox}[Procedimiento]
Completar con los detalles propios de nuestra sesión: desviación inicial
aproximada en la oscilación libre, número de máximos registrados en cada
caso, tiempo de espera hasta el estacionario y posiciones del potenciómetro
fino utilizadas.
\end{pendientebox}

\subsection{Análisis de picos}\label{sec:picos}
Para extraer los máximos de cada registro usamos la función \emph{Análisis de
pico} del menú \emph{Análisis} (\cref{fig:picos}). Su ventana tiene cuatro
elementos:
\begin{itemize}
  \item \emph{Analizar}: elige la curva sobre la que se buscan los picos; con
        un único sensor, la del \emph{eje izquierdo}.
  \item \emph{Tolerancia}: una barra deslizante que decide qué se considera un
        pico. Hacia \emph{menor} detecta también oscilaciones pequeñas; hacia
        \emph{mayor}, solo las más pronunciadas. La ajustamos hasta que el
        programa marca un único pico por oscilación, incluidas las más débiles
        del final del registro.
  \item \emph{Calcular}: genera una tabla con una fila por pico y cinco
        columnas: los instantes de \emph{entrada}, \emph{máximo} y
        \emph{salida} del pico, su \emph{altura} en grados y su \emph{área}.
        De ella tomamos las columnas \emph{máximo}, que da $t_n$, y
        \emph{altura}, que da $\varphi_n$.
  \item \emph{Mostrar resultados}: rotula cada pico sobre la gráfica con su
        instante y su altura, lo que permite comprobar a simple vista que
        ningún máximo se ha perdido ni se ha contado dos veces. Con
        \emph{Copiar} trasladamos la tabla a la hoja de cálculo.
\end{itemize}
Los instantes de la tabla aparecen en múltiplos de \SI{0,1}{\second}, que es la
separación entre dos lecturas consecutivas del sensor; este intervalo de
muestreo fijará la incertidumbre de los tiempos (\cref{sec:incertidumbres}).
Si al final de un registro aparecen picos ajenos a la oscilación, por ejemplo
al detener el anillo con la mano, los descartamos antes del ajuste.

\begin{figure}[!htbp]
\centering
\begin{subfigure}[b]{\textwidth}\centering
  \foto[0.8\textwidth]{figuras/picos-debil.jpg}
  \caption{Amortiguamiento débil.}
\end{subfigure}\\[8pt]
\begin{subfigure}[b]{0.49\textwidth}\centering
  \foto{figuras/picos-medio.jpg}
  \caption{Amortiguamiento intermedio.}
\end{subfigure}\hfill
\begin{subfigure}[b]{0.49\textwidth}\centering
  \foto{figuras/picos-fuerte.jpg}
  \caption{Amortiguamiento fuerte.}
\end{subfigure}
\caption{Análisis de picos sobre tres oscilaciones libres, ordenadas de menor a
mayor amortiguamiento. Las flechas señalan los máximos detectados, rotulados
con su instante y su altura (\emph{Mostrar resultados}), y la ventana
\emph{Análisis de picos} recoge para cada uno los instantes de entrada, máximo
y salida, la altura y el área. Al aumentar el amortiguamiento, la amplitud cae
en menos ciclos: en~(c), en apenas siete máximos pasa de
\SI{77}{\degree} a \SI{12}{\degree}, mientras que en~(a) hacen falta casi
veinte para una caída semejante.}
\label{fig:picos}
\end{figure}

\begin{metodo}[Periodo a partir de los máximos]\label{met:periodo}
Con $N$ máximos consecutivos en los instantes $t_1<t_2<\dots<t_N$ hay $N-1$
periodos completos, de modo que
\begin{equation}\label{eq:periodo}
  T=\frac{t_N-t_1}{N-1} .
\end{equation}
Es la media de los $N-1$ intervalos entre máximos, pero solo intervienen los
dos instantes extremos, y su error se reparte entre todos los periodos: cuantos
más ciclos abarque el registro, más fino es el periodo. Lo aplicamos a la
oscilación libre sin electroimán, para obtener $\wa=2\pi/T$ y de ahí $\wz$
(\cref{obs:w0}), y a cada registro de la oscilación forzada, para obtener
$\Omega=2\pi/T$.
\end{metodo}

\begin{metodo}[Coeficiente de amortiguamiento]\label{met:gamma}
Ajustamos por mínimos cuadrados la recta $\ln\lvert\varphi_n\rvert=a+b\,t_n$ a
los máximos de cada registro y tomamos $\gamma=-b$. Preferimos esta forma a la
de la derecha de la \cref{eq:linealizacion}, con la recta forzada a pasar por
el origen, porque en aquella el primer máximo entra en todos los puntos a
través del cociente $\varphi_0/\varphi_n$: su error se arrastra a toda la serie
y los puntos dejan de ser independientes. Con la ordenada libre, todos los
máximos pesan por igual, y la compatibilidad de $e^{a}$ con $\lvert\varphi_0\rvert$
sirve de comprobación adicional.
\end{metodo}

\begin{metodo}[Amplitud estacionaria]\label{met:beta}
En cada registro de la oscilación forzada, una vez alcanzado el estacionario,
tomamos como amplitud la media de las alturas de los $N$ máximos,
$\beta=\mean{\varphi}_{\text{máx}}$, y como incertidumbre de tipo~A su
desviación típica de la media, $s/\sqrt N$. Si la serie de alturas muestra una
tendencia, creciente o decreciente, es señal de que el transitorio no se había
extinguido y el registro debe descartarse o recortarse.
\end{metodo}

\subsection{Tratamiento de incertidumbres}\label{sec:incertidumbres}
Seguimos la \emph{Guide to the expression of uncertainty in
measurement}~\cite{GUM}, tal como se recoge en las notas del
laboratorio~\cite{NUM-ARG}. Para una magnitud $y=f(x_1,\dots,x_N)$ con
entradas no correlacionadas, la incertidumbre típica se propaga como
\begin{equation}\label{eq:propagacion}
  u^{2}(y)=\sum_{k=1}^{N}\left(\frac{\partial f}{\partial x_k}\right)^{\!2}u^{2}(x_k),
\end{equation}
y la incertidumbre ampliada al \SI{95}{\percent} es $U(y)=t_{95,\nu}\,u(y)$, con
$t_{95,\nu}$ el factor de cobertura de Student para $\nu$ grados de
libertad~\cite{NUM-ARG}; el resultado se expresa como $y\pm U(y)$. Las
incertidumbres se dan con una sola cifra significativa, salvo cuando la
primera es un 1, en cuyo caso se conservan dos, y el valor se redondea a la
misma posición decimal~\cite{ERR-MP}. Dos valores $a$ y $b$ son compatibles
al \SI{95}{\percent} si $\lvert a-b\rvert\leq\sqrt{U^{2}(a)+U^{2}(b)}$.

\subsubsection{Instantes de los máximos}
El sensor se lee cada $\Delta=\SI{0,1}{\second}$, y el máximo verdadero puede
caer en cualquier punto del intervalo de $\pm\Delta/2$ que rodea a la lectura
más alta. Tratándolo como una distribución rectangular de semianchura
$\Delta/2$~\cite{GUM},
\begin{equation}\label{eq:ut}
  u(t_n)=\frac{\Delta/2}{\sqrt3}=\frac{\Delta}{2\sqrt3}\approx\SI{0,03}{\second}.
\end{equation}

\subsubsection{Amplitudes de los máximos}
Por el mismo motivo, la lectura más alta puede no coincidir con la cima de la
oscilación. Si la cima queda a un tiempo $\tau$ de la lectura, con
$\lvert\tau\rvert\leq\Delta/2$, cerca del máximo
$\varphi\simeq\varphi_n\cos(\wa\tau)$, y la altura registrada se queda corta, a
lo sumo, en
\begin{equation}\label{eq:sesgo}
  \frac{\varphi_n-\varphi_{\text{reg}}}{\varphi_n}\leq1-\cos\frac{\wa\Delta}{2}
  \simeq\frac{\wa^{2}\Delta^{2}}{8} .
\end{equation}
Es un error \emph{relativo} y siempre del mismo signo. A ello se suma la
resolución del sensor de ángulo.
\pend{Fijar $u(\varphi_n)$ con el valor de $\wa$ medido y la resolución del
sensor.}

\subsubsection{Periodo y frecuencias}
De la \cref{eq:periodo}, por la \cref{eq:propagacion},
\begin{equation}\label{eq:uT}
  u(T)=\frac{\sqrt2}{N-1}\,u(t_n),\qquad
  u(\wa)=\frac{2\pi}{T^{2}}\,u(T),\qquad
  u(\Omega)=\frac{2\pi}{T^{2}}\,u(T),
\end{equation}
y, para la frecuencia natural corregida, $\wz=\sqrt{\wa^{2}+\gamma^{2}}$,
\begin{equation}\label{eq:uw0}
  u(\wz)=\frac{1}{\wz}\sqrt{\wa^{2}u^{2}(\wa)+\gamma^{2}u^{2}(\gamma)} .
\end{equation}

\subsubsection{Coeficiente de amortiguamiento}
La incertidumbre de $\gamma$ es la de la pendiente del ajuste lineal,
$u(\gamma)=u(b)$, con $\nu=N-2$ grados de libertad. Cada punto lleva una
incertidumbre $u(\ln\lvert\varphi_n\rvert)=u(\varphi_n)/\lvert\varphi_n\rvert$,
que, si $u(\varphi_n)$ es proporcional a $\varphi_n$, es la misma para todos;
en ese caso el ajuste ordinario y el ponderado coinciden.
\pend{Indicar si el ajuste es ponderado y con qué pesos.}

\subsubsection{Frecuencia de resonancia esperada}
De la \cref{eq:resonancia},
\begin{equation}\label{eq:uWR}
  u(\WR)=\frac{1}{\WR}\sqrt{\wz^{2}u^{2}(\wz)+4\gamma^{2}u^{2}(\gamma)} .
\end{equation}

\ornamento

% =========================================================================
\newpage
\seccion{Análisis de datos y resultados}{Lo que dicen las medidas}\label{sec:analisis}
% =========================================================================
\subsection{Frecuencia natural}\label{sec:res-w0}
\begin{pendientebox}[Cuestión 1]
Tabla con los instantes de los máximos de la oscilación libre sin
electroimán. Periodo $T$ por la \cref{eq:periodo}, frecuencia $\wa$ y
frecuencia natural $\wz$ corregida con $\gamma_0$ (\cref{obs:w0}), con sus
incertidumbres típica y ampliada. Comprobación opcional: ajuste lineal de
$t_n$ frente a $n$, cuya pendiente es $T$.
\end{pendientebox}

\subsection{Coeficientes de amortiguamiento}\label{sec:res-gamma}
\begin{pendientebox}[Cuestión 2]
Tabla con $t_n$, $\varphi_n$ y $\ln\lvert\varphi_n\rvert$ para los tres
amortiguamientos (sin puente, puente en~4 y en~6), señalando los puntos
excluidos del ajuste y el motivo. Gráfica semilogarítmica con las tres rectas
de ajuste (\cref{met:gamma}). Parámetros del ajuste con sus incertidumbres y
coeficiente de determinación. Comprobar la curvatura de la
\cref{obs:coulomb} y comparar $\gamma_0$, $\gamma_4$ y $\gamma_6$ entre sí.
\end{pendientebox}

\begin{figure}[!htbp]
\centering
\huecofigura{Gráfica: $\ln\lvert\varphi_n\rvert$ frente a $t_n$ para los tres
amortiguamientos, con sus rectas de ajuste.}
\caption{\pendiente{Pie de figura.}}
\label{fig:ajuste-gamma}
\end{figure}

\subsection{Curvas de resonancia}\label{sec:res-curvas}
\begin{pendientebox}[Cuestión 3]
Tablas con $T$, $\Omega$ y $\beta$ para las siete frecuencias de forzamiento y
los tres amortiguamientos (\cref{met:periodo,met:beta}). Gráfica de $\beta$
frente a $\Omega$ con las tres series y sus barras de incertidumbre.
\end{pendientebox}

\begin{figure}[!htbp]
\centering
\huecofigura[6.2cm]{Gráfica: amplitud estacionaria $\beta$ frente a la
frecuencia de forzamiento $\Omega$ para los tres amortiguamientos, con las
frecuencias de resonancia esperadas $\WR$ señaladas.}
\caption{\pendiente{Pie de figura.}}
\label{fig:curvas}
\end{figure}

\subsection{Frecuencias de resonancia}\label{sec:res-WR}
\begin{pendientebox}[Cuestión 4]
Valores esperados de $\WR$ para cada amortiguamiento a partir de $\wz$ y
$\gamma$ (\cref{eq:resonancia,eq:uWR}), indicados en la gráfica anterior.
Comparación con la frecuencia a la que se observa el máximo de cada curva,
con el criterio de compatibilidad de la \cref{sec:incertidumbres}.
\end{pendientebox}

\subsection{Discusión}\label{sec:discusion-res}
\begin{pendientebox}[Cuestión 5]
Cómo cambian la altura y la anchura del pico con el amortiguamiento
(\cref{sec:resonancia}) y si las razones entre alturas son coherentes con las
razones entre los $\gamma$ medidos. Por qué las tres $\WR$ apenas se
distinguen entre sí. Posibles fuentes de discrepancia: transitorio no
extinguido, amplitudes fuera del régimen lineal del muelle, rozamiento añadido
por la biela al conectar el motor, número de frecuencias medidas cerca del
pico.
\end{pendientebox}

\ornamento

% =========================================================================
\newpage
\seccion{Discusión y conclusiones}{Lo que aprendemos}\label{sec:conclusiones}
% =========================================================================
\begin{pendientebox}[Conclusiones]
Una conclusión por objetivo (frecuencia natural, amortiguamiento, curvas de
resonancia y frecuencias de resonancia), con los valores finales y su
incertidumbre ampliada, seguida de las principales fuentes de error y de las
mejoras que proponemos.
\end{pendientebox}

% =========================================================================
\newpage
%                              BIBLIOGRAFÍA
% =========================================================================
\nocite{CRAWFORD,TIPLER,TAYLOR}
\printbibliography[title={Referencias}]

\end{document}
